\documentclass[11pt,a4paper]{article}
\usepackage[T1]{fontenc}
\usepackage{isabelle,isabellesym}
\usepackage{amssymb}
\usepackage{amsmath}
\usepackage{pdfsetup}

\urlstyle{rm}
\isabellestyle{it}

\begin{document}

\title{The Steiner Deltoid as the Tangent Envelope of Wallace--Simson Lines in Isabelle/HOL}
\author{Arthur Freitas Ramos \and David Barros Hulak \and Ruy J. G. B. de Queiroz}
\maketitle

\begin{abstract}
We formalize the classical Steiner deltoid construction in Isabelle/HOL using
only the published Archive of Formal Proofs entry for the Wallace--Simson line.
In unit-circumcircle coordinates, the deltoid has an explicit complex
parametrisation.  We prove its derivative formula, characterize the stationary
parameters by \(m^3=abc\), and certify them as ordinary cusps using the
second and third derivatives.  We then prove an equality of sets between the
appropriate Wallace--Simson line and the tangent line, for every parameter,
including the cusp cases.

The development uses complex coordinates and handles a vertex parameter by
cycling to the alternate pair of feet.
\end{abstract}

\tableofcontents

\section{Introduction}

The Steiner deltoid is the envelope of the Wallace--Simson lines of a fixed
triangle as the point on its circumcircle varies.  This classical relationship
was treated early by Butchart~\cite{butchart1939} and given a later elementary
proof by de Guzm{\'a}n~\cite{deguzman2001}; computational treatments of
envelopes include~\cite{botana2004}.  The construction is especially
transparent after placing the circumcentre at the origin and scaling the
circumcircle to the unit circle.  If the vertex coordinates are
$a,b,c$ and the moving point is $m$, the deltoid parametrisation is
\[
  d(m) = \frac{a+b+c}{2} + m + \frac{abc\,\overline m^{\,2}}{2}.
\]

The underlying Wallace--Simson theorem is the classical line theorem for the
three perpendicular feet~\cite{wallace1798,coxeter1967}; its Isabelle/HOL
formalization is available as the published AFP entry~\cite{simson-afp}.

This entry deliberately depends only on the published AFP entry
\emph{The Wallace--Simson Line Theorem in Isabelle/HOL}.  In particular, the
perpendicular feet, circle identities, and line interface are imported from
that entry rather than copied from a local development.

\section{Formal setup}

The definition \isa{steiner\_deltoid\_point} gives the displayed
parametrisation, while \isa{steiner\_deltoid} collects its values as $m$ runs
over the unit circle.  For differential statements, the theory introduces
\isa{unit\_circle\_param} and the real curve
\isa{steiner\_deltoid\_curve}.  Its derivative is
\[
  \gamma'(t)=i\,N(e^{it}),\qquad
  N(m)=m-abc\,\overline m^{\,2}.
\]
The tangent direction is defined piecewise: it is \(iN(m)\) when $N(m)\ne0$,
and $m$ when $N(m)=0$.  At a stationary parameter the second derivative is
$-3m$, so this chosen direction is the same tangent direction.
The proof uses the closed form for a perpendicular foot and the published
identity for the difference of two feet on a common circle.

\section{Main result}

The derivative theorem gives
\(\gamma'(t)=iN(e^{it})\).  On the unit circle, the normal vanishes exactly when
\(m^3=abc\), and
\isa{steiner\_deltoid\_stationary\_parameter\_is\_ordinary\_cusp} proves that
the stationary parameters have nonzero, real-linearly independent second and
third derivatives.  The development separately proves that these two
derivatives are perpendicular, a stronger property in this particular curve.
Thus the formal cusp claim uses the standard independence criterion rather
than merely identifying zeros of the first derivative.

The central theorem
\isa{simson\_line\_eq\_steiner\_deltoid\_tangent\_at\_all} states the result as
an equality of sets.  If \(m\ne b\), the line through the \(AB\) and \(BC\)
feet is exactly the tangent at \(d(m)\); if \(m=b\), the line through the
\(BC\) and \(CA\) feet is exactly that tangent.  The proof has separate
regular and stationary branches.  In the regular branch the Simson direction
is perpendicular to \(N(m)\); in the stationary branch the real ratio of the
two feet differences shows that the line has direction \(m\), the direction
of the first nonzero derivative.  The parameter-level theorem
\isa{simson\_line\_eq\_steiner\_deltoid\_tangent\_all} transports this equality
to the real parameter \(t\).

At a stationary parameter, the theorem
\isa{steiner\_deltoid\_cusp\_tangent\_is\_second\_derivative} explicitly
identifies the piecewise-defined cusp tangent with the line in the direction
of \(\gamma''(t)=-3m\).  The older foot-membership results are named
\isa{simson\_line\_is\_tangent\_to\_steiner\_deltoid\_at\_regular} and
\isa{simson\_line\_is\_tangent\_to\_steiner\_deltoid\_regular} to make their
non-stationary hypotheses visible in the public API.

\input{session}

\bibliographystyle{abbrv}
\bibliography{root}

\end{document}
